% !TeX root = ../main.tex
\section{Conclusion}
\label{sec:conclusion}

We presented a reproducible study of deterministic QUIC metadata obfuscation against
deep learning traffic classifiers on CESNET-QUIC22 backbone data. Seven obfuscation
policies expose a structured privacy--cost landscape: timing jitter buys privacy at
latency cost with zero bandwidth expansion; padding buys stronger privacy at heavy
bandwidth cost. A masked Transformer attacker achieves 74.4\% macro F1 on clean held-out
flows; \emph{jitter\_low} preserves 72.7\% macro F1 at 11\,ms overhead---a practical
proxy configuration supported by bootstrap and McNemar tests. Transformers outperform
CNN-BiLSTMs under jitter, while MTU padding defeats both. Our Pareto analysis and open
pipeline aim to support evidence-based deployment of metadata defenses in QUIC-heavy
networks.

%==============================================================================
